% Figure: normal strain (extension) and shear strain (angle change) of a 2D
% element, illustrating epsilon_xx as fractional length change and gamma_xy as
% total angle change.
\begin{figure}[htbp]
\centering
\begin{tikzpicture}[scale=1.0, >=latex]
% ==== Left: normal strain ====
\begin{scope}[shift={(0,0)}]
  \draw[enggray, thick] (0,0) rectangle (2,2);
  \draw[engblue, very thick, dashed] (0,0) rectangle (2.6,1.8);
  \draw[<->, engblue] (0,-0.35) -- (2.6,-0.35);
  \node[engblue, font=\scriptsize, anchor=north] at (1.3,-0.4) {$L+\Delta L$};
  \draw[<->, enggray] (0,2.3) -- (2,2.3);
  \node[enggray, font=\scriptsize, anchor=south] at (1,2.3) {$L$};
  \node[font=\small, anchor=north] at (1.3,-1.0) {normal strain $\epsilon_{xx}=\Delta L/L$};
\end{scope}
% ==== Right: shear strain ====
\begin{scope}[shift={(6,0)}]
  \draw[enggray, thick] (0,0) rectangle (2,2);
  % sheared parallelogram
  \draw[engred, very thick, dashed] (0,0) -- (0.6,2) -- (2.6,2) -- (2,0) -- cycle;
  % angle arc
  \draw[engred, ->] (0,1) arc [start angle=90, end angle=72, radius=1];
  \node[engred, font=\scriptsize] at (0.55,1.2) {$\gamma_{xy}$};
  \node[font=\small, anchor=north] at (1.3,-1.0) {shear strain $\gamma_{xy}=2\epsilon_{xy}$};
\end{scope}
\end{tikzpicture}
\caption[Normal strain and shear strain]{The two kinds of strain on a
two-dimensional element. Normal strain $\epsilon_{xx}=\Delta L/L$ is the
fractional change in length of a line originally along the $x$-axis
(left). Engineering shear strain $\gamma_{xy}$ is the total change in the
right angle between two lines originally along the $x$- and $y$-axes
(right); the tensor shear strain $\epsilon_{xy}$ is half of it,
$\gamma_{xy}=2\epsilon_{xy}$. The factor of two is the most common source
of error in mechanics-of-materials calculation: Hooke's law for shear is
written $\tau=G\gamma$ with the engineering strain, not $2G\epsilon_{xy}$.}
\label{fig:vol03ch03:strain-deformation}
\end{figure}
